\begin{titlepage}
    \centering
    \vspace*{1.0cm}
    
    {\scshape\Large Universidad / Instituto Tecnológico \par}
    \vspace{0.3cm}
    {\scshape\large Asignatura: Internet de las Cosas (IoT) / Sistemas Embebidos \par}
    \vspace{0.2cm}
    {\large Unidad 2: Comunicación Bidireccional IoT \par}
    
    \vspace{2.0cm}
    
    \begin{tcolorbox}[colback=primaryBlue,colframe=primaryBlue,arc=2mm,center,width=0.95\textwidth]
        \centering\color{white}
        \vspace{0.3cm}
        {\LARGE\bfseries REPORTE TÉCNICO \par}
        \vspace{0.25cm}
        {\Large Comunicación Bidireccional entre ESP32 y Node-RED mediante Protocolo MQTT \par}
        \vspace{0.2cm}
        {\small Telemetría de Temperatura/Humedad, Control de Actuador LED y Supervisión LWT \par}
        \vspace{0.3cm}
    \end{tcolorbox}
    
    \vspace{2.5cm}
    
    \begin{minipage}[t]{0.85\textwidth}
        \centering \large
        \textbf{Autores:}\\
        \vspace{0.2cm}
        {\Large Arturo Heredia \quad \textbar \quad Diego López}
    \end{minipage}
    
    \vspace{2.5cm}
    
    \begin{infobox}[Código Fuente y Recursos]
        \centering\small
        Repositorio GitHub: \textbf{\url{https://github.com/arturoherediam/esp32-monitoreo-ambiental-iot}}
    \end{infobox}
    
    \vfill
    
    {\large Octubre de 2026 \par}
\end{titlepage}
